\chapter{INTRODUCTION}

\ChapterSection{Background}

Agriculture is not only an economic sector in Nepal; it is a daily livelihood system for a large
part of the population. The Government of Nepal publishes national agricultural statistics to
track crop production, livestock, land use, and the role of agriculture in the national economy
\parencite{moaldAgricultureStats2080}. The National Statistics Office also conducted the National
Sample Census of Agriculture 2021/22 across all districts and local levels to provide benchmark
data on agricultural holdings, land use, crop areas, livestock, irrigation, machinery, labour,
and agricultural inputs \parencite{nsoAgricultureCensus2021}. These official efforts show that
agriculture remains a planning, livelihood, and food-security priority for Nepal.

Crop disease diagnosis sits at the practical centre of this problem. A disease symptom often
appears first as a visual change on a leaf, fruit, panicle, ear, or other plant structure. At
that moment, the farmer needs timely guidance: whether the symptom is serious, whether it is
likely to spread, what crop part is affected, and whether expert consultation is needed. If this
decision is delayed or made incorrectly, the effect can move beyond one plant and affect crop
yield, input cost, and household income. FAO reports that plant pests and diseases can cause up
to 40 percent annual crop-production losses globally, which underlines the seriousness of
plant-health problems even though the exact loss depends on crop, location, and season
\parencite{faoPlantHealthLosses}.

\ChapterSubsection{Crop disease diagnosis challenge}

Traditional disease diagnosis depends on farmer experience, local agro-vet advice, extension
workers, or laboratory confirmation. These channels are important, but they are not always
available at the exact time a farmer observes a symptom in the field. Nepal's geography, rural
settlement patterns, variable connectivity, and uneven access to specialist services make this
problem more difficult. A farmer may have a phone and a crop image, but still lack immediate
access to a trained plant-health expert.

The diagnosis challenge is also technical. Field images are not controlled laboratory samples.
They may include unclear lighting, cluttered background, multiple plant parts, partial symptoms,
low image quality, or mixed disease stages. A reliable digital system must therefore do more than
classify a clean image. It must support image capture, model inference, structured interpretation,
advisory output, local record keeping, and uncertainty handling.

\ChapterSubsection{Computer vision and mobile advisory context}

Computer vision can support agricultural diagnosis by extracting information from crop images.
Object detection can identify relevant plant structures, while vision-language models can
produce disease-oriented interpretation from visual input. Mobile technology is important because
the smartphone is the most realistic interface for many field-level users. A diagnosis-support
system therefore becomes useful only when the model layer is connected to a mobile workflow and a
backend service that can store records, manage users, validate requests, and return structured
advisories.

This direction is aligned with Nepal's movement toward digital agriculture. FAO reports that
Nepal's digital agriculture strategy work emphasizes digital infrastructure, digital skills,
public-private coordination, advisory services, and rural livelihood improvement
\parencite{faoNepalDigitalAgriculture}. Krishi Vaidya is positioned within this direction as a
student-built integrated prototype for crop-health support.

\ChapterSubsection{Krishi Vaidya project context}

Krishi Vaidya is an integrated final-year B.Sc. CSIT project composed of four major technical
subsystems. The YOLO subsystem performs plant-part object detection. The VLM subsystem performs
image-based disease interpretation and structured output generation. The backend subsystem
provides authentication, records, scan handling, diagnosis orchestration, advisory generation,
validation, and persistence. The mobile application provides the farmer-facing interface for
image capture, scan submission, diagnosis viewing, crop records, local storage, and offline-aware
workflow.

The project is therefore not a single machine-learning script. It is a full diagnosis-support
prototype where trained model components are connected to application software. This manuscript
documents the design, implementation, results, limitations, and remaining validation needs of
that integrated system.

\ChapterSection{Problem Statement}

Farmers need timely, understandable, and accessible support for crop disease diagnosis, but
expert diagnosis is not always available at the field level. Existing digital approaches often
focus on isolated image classification or model accuracy without fully addressing the workflow
required by a practical advisory system. A useful crop-health support system must connect image
capture, object localization, disease interpretation, backend validation, advisory generation,
mobile display, local records, and failure handling.

The problem addressed by this project is therefore:

\begin{quote}
How can an integrated mobile, backend, object-detection, and vision-language system be designed
and implemented to support image-based crop-health diagnosis and farmer-facing advisory output
for the Nepalese agricultural context?
\end{quote}

\ChapterSubsection{Visual detection problem}

Crop images contain different plant structures, and disease-related reasoning often depends on
which part of the plant is visible. A system that ignores plant-part localization may treat a
complex field image as a simple classification input. Krishi Vaidya addresses this by training a
YOLO object detector on agricultural plant-part classes such as leaves, panicles, grain clusters,
fruits, ears, and vegetable heads.

\ChapterSubsection{Disease interpretation problem}

Disease diagnosis requires more than detecting a plant object. The system must interpret visual
evidence and produce structured information that can be converted into an advisory response.
Krishi Vaidya addresses this through a VLM pipeline trained and evaluated for disease-oriented
output. The project also records the current weakness of this component: prediction quality and
parseability require further improvement before the VLM can be treated as a robust diagnostic
authority.

\ChapterSubsection{Mobile access and advisory workflow problem}

Even a useful model is incomplete if farmers cannot access it through a practical interface. The
mobile app must support image capture, diagnosis requests, offline records, synchronization,
crop records, and understandable result presentation. The backend must protect the API, validate
requests, call the diagnosis service, store records, and return consistent responses. Krishi
Vaidya addresses this by implementing both the mobile and backend layers rather than limiting the
project to model training.

\ChapterSection{Objectives}

\ChapterSubsection{General objective}

The general objective of Krishi Vaidya is to design and implement an integrated crop-health
diagnosis-support system that uses object detection, vision-language modelling, backend services,
and a mobile application to assist farmers with image-based crop disease interpretation and
advisory output.

\ChapterSubsection{Specific objectives}

\begin{enumerate}[label=\arabic*.]
\item To prepare and train a YOLO-based object-detection model for agricultural plant-part
localization.
\item To prepare and fine-tune a VLM-based disease interpretation pipeline using structured
image-text data.
\item To implement a backend API for user authentication, crop records, scan handling, diagnosis
orchestration, advisory generation, validation, and persistence.
\item To implement a mobile application for image-based scanning, diagnosis display, local
records, offline-aware storage, synchronization, and farmer-facing output.
\item To integrate the mobile app, backend, and model services into a documented diagnosis
workflow.
\item To evaluate the implemented components using available model metrics, validation outputs,
source-backed tests, generated diagrams, and documented evidence gaps.
\end{enumerate}

\ChapterSection{Scope}

The scope of this project includes the design, implementation, and manuscript documentation of
the Krishi Vaidya prototype. The model scope includes YOLO object detection for plant-part
localization and VLM-based disease interpretation. The backend scope includes the implemented
Express/Mongoose service structure, scan and diagnosis workflow, advisory response handling,
security middleware, validation, and tests available in the repository. The mobile scope includes
the implemented React Native/Expo app structure, scan workflow, local SQLite persistence,
offline-aware behavior, synchronization structure, crop and task records, and farmer-facing
diagnosis screens.

The scope does not include claiming a certified medical, agronomic, or government-approved
diagnosis system. The project is a research and engineering prototype. It is suitable for
demonstrating integrated software and model development, but it requires further expert review,
field testing, and deployment hardening before real-world advisory use.

\ChapterSection{Limitations}

The main limitations of the project are as follows:

\begin{enumerate}[label=\arabic*.]
\item The YOLO dataset is large and curated, but class distribution is imbalanced; some classes
have much fewer annotations than others.
\item The VLM pipeline is implemented and validated, but the current prediction metrics and
parse-error count show that structured diagnostic reliability must be improved.
\item The backend has scan-focused test evidence, but broader endpoint coverage for all domains
is still required.
\item The mobile app implementation is source-verified, but final screenshots, physical-device
testing, runtime latency, and offline synchronization demonstrations still need to be captured.
\item The integration workflow is source-verified, but live mobile-to-backend-to-VLM
request-response evidence remains a final validation requirement.
\item The advisory output has not yet been validated by certified agricultural experts and should
not be used as a replacement for expert consultation.
\end{enumerate}

\ChapterSection{Methodology Overview}

The project was implemented using an evidence-first workflow. First, the YOLO dataset was
constructed, curated, trained, validated, and evaluated using detection metrics, threshold
analysis, quantization comparison, and error analysis. Second, the VLM dataset was prepared using
structured image-text records, then fine-tuned using LoRA and quantized training methods, with
evaluation focused on dataset validity, training behavior, resource use, prediction metrics, and
parseability. Third, the backend was implemented as a layered API with routes, controllers,
services, repositories, models, middleware, and tests. Fourth, the mobile app was implemented as
a React Native/Expo application with routing, feature modules, local persistence, scan workflow,
backend communication, and offline-aware behavior. Finally, the integration workflow was mapped
across the mobile, backend, VLM, and advisory layers, and the testing/validation evidence was
documented.

\ChapterSection{Report Organization}

This report is organized into five chapters. Chapter 1 introduces the project context, problem,
objectives, scope, limitations, and methodology overview. Chapter 2 reviews the literature and
technical background related to digital agriculture, crop disease diagnosis, YOLO object
detection, vision-language models, mobile systems, backend systems, and the research gap.
Chapter 3 presents the materials and methods used to build the YOLO model, VLM pipeline,
backend, mobile app, integration workflow, and validation structure. Chapter 4 presents results
and discussion, including model metrics, implementation outcomes, integration interpretation,
limitations, and evidence gaps. Chapter 5 concludes the work and provides recommendations for
future improvement. The appendix provides supplementary implementation and evidence material.
